\pdfbookmark[1]{Extract from three years of daily observations of magnetic declination at Göttingen.}{dreijaehrige-declination}
\begin{center}\large\textit{Extract from three years of daily observations of magnetic declination at Göttingen.}\end{center}
\begin{center}\small Results of the Observations of the Magnetic Association.\quad 1836.\quad III.\end{center}

Amid the incessant succession of smaller and larger fluctuations in magnetic declination, which we call irregular insofar as their occurrence follows no rule of timing, the only way to separate out what is regular is to make a great many observations according to a definite plan, to continue them with persistent consistency for a long time, and to derive means by combining them suitably. In those means the influence of the anomalies that always affect individual observations disappears as far as possible. During the morning hours the declination usually increases in our region, but by a great deal one day and by little the next. Sometimes (though rarely) one even observes a smaller declination at the hour when it is usually greatest than in the early hours of the same day. The cause of the morning increase may well be active every day, but its effect is sometimes increased, sometimes diminished, and sometimes completely obscured by other forces that intervene irregularly. How much effect the regular cause actually has, and how its effect differs from one season to another, can be
\SourceLeaf{563}ascertained only from means taken over very many days. A single day or a few days do not suffice. The case is similar for the gradual changes that progress in the same direction for at least a very long time. We call these secular changes because it takes a long series of years for them to add up to many degrees. Individual observations only a few years apart cannot give us any reliable information about this yet, even if they are made on the same day of the month and at the same hour. But means obtained consistently already allow us after a few years to anticipate what could otherwise be established with some approximation only after several decades.

With this in mind, from the outset I included among the observations to be made at our magnetic observatory the daily determination of absolute declination, always at the same hours. Work of this kind derives its value from being continued for a long time without interruption. To make it easier to count on being able to do so, I preferred to start with a limited plan and avoid trying to cover too much at once. Hence only two determinations are made daily, at 8 in the morning and 1 in the afternoon, mean time. Besides fitting most easily around other duties, these hours had the advantage that, when magnetic movements follow their regular course, the position of the needle at 1 in the afternoon is always close to the maximum declination, while at 8 in the morning it is close to the minimum for most of the year. Observing at fixed hours of \textit{true} solar time would in itself have been somewhat more natural. Here, however, where what mattered above all was consistent application of the same principle, the far greater convenience of an arrangement based on mean time had to prevail.

These regular records began on the first of January 1834. The first two and a half months have been excluded from the extract below, however. During that time the suspension thread often had to be wound up, which changed the torsion zero point in ways that received too little attention at first. From 17 March onward a stronger suspension thread (of two hundred strands) was used, after its torsion zero point had first been adjusted precisely. Whenever a change was subsequently made to this thread or to any other circumstance affecting the
\SourceLeaf{564}reduction elements, the necessary corrections or changes to the reduction elements were made on each occasion. In the first months several sufficiently practised observers shared the observations with me. From 1 October 1834 onward they have been made regularly by Dr.~\textsc{Goldschmidt}, with other skilled observers taking his place only when he is prevented from doing so.

I have already reported the monthly means of these determinations through January 1835 in the \textit{Göttingische gelehrte Anzeigen} for 1834 and 1835 [pp.~519 and 528 of this volume]. The means for three full years now follow.

\Needspace{31\baselineskip}
\begin{center}\textit{Mean westward magnetic declination at Göttingen.}\end{center}
\begin{center}\small
\begin{tabular}{rl|rrr|rrr}
&&\multicolumn{3}{c|}{8 a.m.}&\multicolumn{3}{c}{1 p.m.}\\\hline
1834&March, second half&$18^0$&$38'$&$16''0$&$18^0$&$46'$&$40''4$\\
&April&&36&6.9&&47&3.8\\
&May&&36&28.2&&47&15.4\\
&June&&37&40.7&&47&59.5\\
&July&&37&57.5&&48&19.0\\
&August&&38&48.1&&49&11.0\\
&September&&36&58.4&&46&32.3\\
&October&&37&18.4&&44&47.2\\
&November&&37&38.4&&43&4.3\\
&December&&37&54.8&&41&32.7\\[4pt]
1835&January&&37&51.5&&42&14.4\\
&February&&37&3.5&&42&29.4\\
&March&&34&47.5&&44&55.2\\
&April&&32&57.7&&46&31.6\\
&May&&32&13.4&&45&17.1\\
&June&&32&56.4&&44&41.3\\
&July&&34&8.0&&44&42.8\\
&August&&34&12.4&&46&56.8\\
&September&&33&21.2&&44&27.6\\
&October&&33&23.0&&43&5.3\\
&November&&36&15.3&&43&49.5\\
&December&&35&25.9&&40&19.1\\[4pt]
1836&January&&35&2.4&&40&34.6\\
&February&&33&26.7&&41&15.2\\
&March&&31&1.4&&43&16.4\\
&April&&26&32.9&&43&42.6\\
&May&&28&0.8&&44&37.2
\end{tabular}
\end{center}
\SourceLeaf{565}\Needspace{15\baselineskip}
\begin{center}\small
\begin{tabular}{rl|rrr|rrr}
&&\multicolumn{3}{c|}{8 a.m.}&\multicolumn{3}{c}{1 p.m.}\\\hline
1836&June&$18^0$&$27'$&$35''1$&$18^0$&$42'$&$52''4$\\
&July&&26&54.2&&42&26.0\\
&August&&25&42.4&&41&45.0\\
&September&&26&14.6&&40&59.6\\
&October&&27&34.0&&40&32.8\\
&November&&29&21.0&&36&54.3\\
&December&&29&13.7&&35&46.8\\[4pt]
1837&January&&27&35.3&&37&46.2\\
&February&&27&35.6&&36&28.3\\
&March&&25&44.2&&39&4.2
\end{tabular}
\end{center}

A few combinations of these observations may now be given.

The difference between morning and afternoon declination has the same sign throughout the mean values. The following table shows how the size of the difference depends on the season.
\begin{center}\small
\begin{tabular}{l||rr|rr|rr||rr}
&\multicolumn{2}{c|}{1834. 1835}&\multicolumn{2}{c|}{1835. 1836}&\multicolumn{2}{c||}{1836. 1837}&\multicolumn{2}{c}{Mean}\\\hline\hline
April&$10'$&$56''9$&$13'$&$33''9$&$17'$&$9''7$&$13'$&$53''5$\\
May&10&47.2&13&3.7&16&36.4&13&29.1\\
June&10&18.8&11&44.9&15&17.3&12&27.0\\
July&10&21.5&10&34.8&15&31.8&12&9.4\\
August&10&22.9&12&44.4&16&2.6&13&3.3\\
September&9&33.9&11&6.4&14&45.0&11&48.4\\
October&7&28.8&9&42.3&12&58.8&10&3.3\\
November&5&25.9&7&34.2&7&33.3&6&51.1\\
December&3&37.9&4&53.2&6&33.1&5&1.4\\
January&4&22.9&5&32.2&10&10.9&6&42.0\\
February&5&25.9&7&48.5&8&52.7&7&22.4\\
March&10&7.7&12&15.0&13&20.0&11&54.2\\\hline\hline
Mean&$8'$&$14''2$&$10'$&$2''8$&$12'$&$54''3$&$10'$&$23''8$
\end{tabular}
\end{center}
We see that the difference was smallest in December, both in the means and in each individual year. This seems quite natural, since changes that vary with the time of day must necessarily be attributed to an effect of the Sun, even though we do not yet know \textit{how} that effect comes about. What may at first seem surprising, on the other hand, is that the much larger differences in the summer months were not greatest around the
\SourceLeaf{566}solstice and were smaller in June and July than in April, May and August. This is especially so since the agreement of all three individual years in this respect gives reason to suppose that it is no accident. It must be borne in mind, however, that in the months nearest the solstice the minimum declination occurs at an earlier hour, so the total increase would be appreciably greater than the movement reckoned from 8 o’clock.

It is also striking that the difference was greater in every month of the second year than in the first, and again greater in the third than in the second. The differences are far too large to be taken as evidence of a secular increase. Rather, if the observations are continued for several years, fluctuations in either direction are to be expected. In any event, we learn from this that one year may stand out from another in the Sun’s effect on terrestrial magnetism, much as a whole summer or a whole winter differs considerably from others in its weather. For that very reason it will take several years of observations to determine the mean values accurately.

It has already been noted above that on individual days the difference between morning and afternoon declination can exceptionally have the opposite sign. The rarity of these exceptions is clear from the fact that only fourteen such cases occurred in three years of observations, or an average of one in 79 days. I list them here, together with the amount by which the declination at 8 in the morning exceeded that at 1 in the afternoon on each occasion.
\begin{center}\small
\begin{tabular}{rlr|rr||rlr|rr}
1834&August&15&$6'$&$8''0$&1835&November&8&$3'$&$42''2$\\
&December&24&3&43.0&&December&8&18&35.6\\
&December&25&0&38.2&1836&January&20&0&46.3\\
&December&26&2&20.3&&July&20&5&8.8\\
1835&January&30&0&23.8&&November&9&11&9.5\\
&February&7&0&32.5&1837&February&13&4&1.0\\
&October&4&0&43.1&&March&14&1&22.6
\end{tabular}
\end{center}
It is quite natural that twelve of these fourteen exceptions fall in the winter months and only two in the summer months. An anomalous movement can more readily outweigh the small
\SourceLeaf{567}regular effect of the Sun in the former than the much larger effect in the latter.

To see how far the secular change can already be discerned from the observations at hand, the monthly means of the first year have been compared with the corresponding means of the second, and those of the second likewise with those of the third. Of the 48 comparisons obtained in this way (the incomplete month of March 1834 is excluded from this and the other combinations), 47 give a decrease. Only one gives an increase, and it is therefore marked with a minus sign in the following table.
\Needspace{17\baselineskip}
\begin{center}\small\setlength{\tabcolsep}{3pt}
\textit{Annual decrease in declination.}\par\smallskip
\begin{tabular}{l||rr|rr||rr|rr||rr}
&\multicolumn{4}{c||}{First year}&\multicolumn{4}{c||}{Second year}&\multicolumn{2}{c}{Mean}\\
&\multicolumn{2}{c|}{8 a.m.}&\multicolumn{2}{c||}{1 p.m.}&\multicolumn{2}{c|}{8 a.m.}&\multicolumn{2}{c||}{1 p.m.}&&\\\hline\hline
April&$3'$&$9''2$&$0'$&$32''2$&$6'$&$24''8$&$2'$&$49''0$&$3'$&$13''8$\\
May&4&14.8&1&58.3&4&12.6&0&39.9&2&46.4\\
June&4&44.3&3&18.2&5&21.3&1&48.9&3&48.1\\
July&3&49.5&3&36.2&7&13.8&2&16.8&4&14.1\\
August&4&35.7&2&14.2&8&30.0&5&11.8&5&7.9\\
September&3&37.2&2&4.7&7&6.6&3&28.0&4&4.1\\
October&3&55.4&1&41.9&5&49.0&2&32.5&3&29.6\\
November&1&23.1&$-0$&45.2&6&54.3&6&55.2&3&36.8\\
December&2&28.9&1&13.6&6&12.2&4&32.3&3&36.7\\
January&2&49.1&1&39.8&7&27.1&2&48.4&3&41.1\\
February&3&36.8&1&14.2&5&51.1&4&46.9&3&52.2\\
March&3&46.1&1&38.8&5&17.2&4&12.2&3&43.6\\\hline\hline
Mean&3&30.8&1&42.2&6&21.7&3&20.2&3&46.2
\end{tabular}
\end{center}
That the comparison of the morning means usually gives a larger decrease here than the comparison of the afternoon means is simply another way of expressing what was observed above, that the daily changes were found to be smaller in the first year than in the second, and smaller in the second than in the third. The difference should therefore be regarded as due to chance rather than real, and if observations are continued for longer, a difference in the opposite direction is also to be expected. Insofar, then, as there is no sufficient reason to prefer one result to the other, all that remains is to take the mean of
\SourceLeaf{568}the two. This mean is $2'36''5$ for the first year and $4'55''9$ for the second, and one might be tempted to take this as proof that the decrease in declination is accelerating. That, however, would be no more than a poor argument for a claim that is true in itself. It is known that the declination, which increased throughout Europe during the preceding century, reached its maximum in the present century and has been falling again ever since. The nature of the case requires this turning point to produce a decrease that is imperceptible at first and becomes gradually stronger. For lack of earlier observations, the year when the turning point occurred for Göttingen cannot be specified with certainty. But from the observations known from other places, that year must be taken to lie considerably further back than those two numbers would imply if they were regarded as pure effects of the slow movement we call secular motion. Likewise, in the light of all other experience, a change as great as $2'19''4$ is quite inadmissible as a regular increase over one year. We therefore regard this difference as largely due to chance as well. For the time being, until more extensive experience is available, the mean $3'46''2$ must accordingly be taken as the annual decrease in declination for 1834---1837.

Since the difference between the declinations at the morning and afternoon hours varies so clearly with the season, the question arises whether only one of the two, or one in particular, or both together take part in a change that depends on the season, and what laws underlie it. To determine these laws, a longer series of years will be even more necessary than it is for the difference in declination alone. In the meantime, however, it is worth seeing what the observations to date can tell us within their limits.

For this purpose the means over twelve months have first been calculated for each of the three years of observation. They are given below.
\begin{center}\small
\begin{tabular}{l|rrr|rrr}
&\multicolumn{3}{c|}{8 a.m.}&\multicolumn{3}{c}{1 p.m.}\\\hline
1834---1835&$18^0$&$37'$&$12''5$&$18^0$&$45'$&$27''0$\\
1835---1836&&33&42.0&&43&44.8\\
1836---1837&&27&20.3&&40&14.6
\end{tabular}
\end{center}

\SourceLeaf{569}These means are to be regarded as applying to the middle day of each year of calculation, the first pair thus applying to 1 October 1834, and so on.

Comparing each month of each year with the corresponding annual mean gives the following differences.
\begin{center}\small
\textit{Declination at 8 a.m.}\par\smallskip
\begin{tabular}{l|rr|rr|rr||rr}
&\multicolumn{2}{c|}{First year}&\multicolumn{2}{c|}{Second year}&\multicolumn{2}{c||}{Third year}&\multicolumn{2}{c}{Mean}\\\hline
April&$-1'$&$5''9$&$-0'$&$44''3$&$-0'$&$47''4$&$-0'$&$52''5$\\
May&$-0$&44.6&$-1$&28.6&$+0$&40.5&$-0$&30.9\\
June&$+0$&27.9&$-0$&45.6&$+0$&14.8&$-0$&1.0\\
July&$+0$&44.7&$+0$&26.0&$-0$&26.1&$+0$&14.9\\
August&$+1$&35.3&$+0$&30.4&$-1$&37.9&$+0$&9.3\\
September&$-0$&14.4&$-0$&20.8&$-1$&5.7&$-0$&33.6\\
October&$+0$&5.6&$-0$&19.0&$+0$&13.7&$-0$&0.1\\
November&$+0$&25.6&$+2$&33.3&$+2$&0.7&$+1$&39.9\\
December&$+0$&42.0&$+1$&43.9&$+1$&53.4&$+1$&26.4\\
January&$+0$&38.7&$+1$&20.4&$+0$&15.0&$+0$&44.7\\
February&$-0$&9.3&$-0$&15.3&$+0$&15.3&$-0$&3.1\\
March&$-2$&25.3&$-2$&40.6&$-1$&36.1&$-2$&14.0
\end{tabular}
\end{center}
\Needspace{16\baselineskip}
\begin{center}\small
\textit{Declination at 1 p.m.}\par\smallskip
\begin{tabular}{l|rr|rr|rr||rr}
&\multicolumn{2}{c|}{First year}&\multicolumn{2}{c|}{Second year}&\multicolumn{2}{c||}{Third year}&\multicolumn{2}{c}{Mean}\\\hline
April&$+1'$&$36''8$&$+2'$&$46''8$&$+3'$&$28''0$&$+2'$&$37''2$\\
May&$+1$&48.4&$+1$&32.3&$+4$&22.6&$+2$&34.4\\
June&$+2$&32.5&$+0$&56.5&$+2$&37.8&$+2$&2.3\\
July&$+2$&52.0&$+0$&58.0&$+2$&11.4&$+2$&0.5\\
August&$+3$&44.0&$+3$&12.0&$+1$&30.4&$+2$&48.8\\
September&$+1$&5.3&$+0$&42.8&$+0$&45.0&$+0$&51.0\\
October&$-0$&39.8&$-0$&39.5&$+0$&18.2&$-0$&20.4\\
November&$-2$&22.7&$+0$&4.7&$-3$&20.3&$-1$&52.8\\
December&$-3$&54.3&$-3$&25.7&$-4$&27.8&$-3$&55.9\\
January&$-3$&12.6&$-3$&10.2&$-2$&28.4&$-2$&57.1\\
February&$-2$&57.6&$-2$&29.6&$-3$&46.3&$-3$&4.5\\
March&$-0$&31.8&$-0$&28.4&$-1$&10.4&$-0$&43.5
\end{tabular}
\end{center}
As means over three years, the numbers in the last column have been freed to some extent, though still very imperfectly, from the influence of irregular anomalies. They are evidently still affected by the secular change. To remove it, its amount between the
\SourceLeaf{570}middle of each month and 1 October must still be applied with a negative sign for the first six months and with a positive sign for the last six. Using the twelve-month value $3'46''2$ determined above as the basis, we obtain the following results.
\begin{center}\small
\begin{tabular}{l|rr|rr||rr}
&\multicolumn{2}{c|}{8 a.m.}&\multicolumn{2}{c||}{1 p.m.}&\multicolumn{2}{c}{Mean}\\\hline
April&$-2'$&$35''6$&$+0'$&$54''2$&$-0'$&$50''7$\\
May&$-1$&55.3&$+1$&10.0&$-0$&22.6\\
June&$-1$&6\phantom{.}6&$+0$&56.7&$-0$&4.9\\
July&$-0$&32.0&$+1$&13.6&$+0$&20.8\\
August&$-0$&18.8&$+2$&20.7&$+1$&0.9\\
September&$-0$&43.0&$+0$&41.6&$-0$&0.7\\
October&$+0$&9.3&$-0$&11.0&$-0$&0.8\\
November&$+2$&8.0&$-1$&24.7&$+0$&21.6\\
December&$+2$&13.3&$-3$&9.0&$-0$&27.8\\
January&$+1$&50.3&$-1$&51.5&$-0$&0.6\\
February&$+1$&21.3&$-1$&40.1&$-0$&9.4\\
March&$-0$&30.9&$+0$&59.6&$+0$&14.3
\end{tabular}
\end{center}
These results already show as much regularity as could be expected from only three years of observations. The first column shows how much the morning declination in each month differs from the mean morning declination. Likewise, the second column gives the difference between the afternoon declination in each month and the mean afternoon declination. It must be remembered that the latter itself is $10'23''8$ greater than the mean morning declination.

It seems remarkable that in all twelve months the morning and afternoon declinations fluctuate \textit{on opposite sides} of their mean values. In the five winter months from October to February the morning declination is greater than its mean and the afternoon declination smaller, so both circumstances contribute \textit{at the same time} to bringing the total difference below its mean value in that season. The exact opposite occurs in the other seven months. Moreover, these opposite fluctuations are on average nearly equal in size. As a result they almost cancel in their mean, given in the last column. In other words, the mean of the magnetic declinations at 8 in the morning and 1 in the afternoon
\SourceLeaf{571}contains, besides the irregular anomalies and the secular decrease, no appreciable fluctuations depending on the season. At least no difference between the summer and winter months emerges with certainty.

The mean value itself, derived from all three years of observations, would be written for 1 October 1835 as follows
\[
=18^0\;37'\;56''9
\]
It goes without saying, of course, that the mean here is taken only over the hours chosen for our observations. The mean over \textit{all} the hours of the day might differ somewhat from it, though probably only slightly. But the investigations so far show amply that nothing can be settled with certainty on that score without very prolonged work.

So far we have dealt only with the monthly means. Printing all the individual observations was deemed unnecessary for the time being, because as long as they come from only one place their interest would lie solely in what they reveal about the irregular jumps back and forth. A systematic combination of the numbers serves this purpose better than merely looking at them. It reduces the extent of the fluctuations to a definite measure and makes it possible to compare the general character of different periods precisely as to whether fluctuations were greater or smaller during them. For brevity, by a fluctuation of magnetic declination I mean here the difference from its value on the preceding day at the same hour. By the mean fluctuation during any period (by analogy with the so-called mean errors of observation) I mean the square root of the mean of the squares of the individual fluctuations. Note that if several periods that are equal or treated as equal are subsequently to be combined into a single period, one may not take the arithmetic mean of their separate mean fluctuations to determine the overall mean. One must first go back to the squares of the separate means, find their arithmetic mean and take its square root. The results of calculating in this way from the three years of
\SourceLeaf{572}observations are given in the following table, in seconds of arc.
\Needspace{19\baselineskip}
\begin{center}\small
\textit{Mean fluctuation of magnetic declination during the three years}\par
1834---1837.\par\smallskip
\begin{tabular}{l|r|r|r||r||r|r|r||r}
&\multicolumn{4}{c||}{8 a.m.}&\multicolumn{4}{c}{1 p.m.}\\\cline{2-9}
&I&II&III&Mean&I&II&III&Mean\\\hline
April&74&126&205&147&129&101&264&180\\
May&192&124&277&207&158&183&210&185\\
June&172&171&199&181&95&151&217&162\\
July&213&243&287&250&119&184&252&193\\
August&264&253&269&262&175&165&307&225\\
September&162&325&207&241&172&143&161&159\\
October&116&296&216&222&182&202&242&210\\
November&79&205&308&218&170&173&126&158\\
December&132&324&71&206&184&206&154&182\\
January&146&274&138&196&174&212&154&181\\
February&116&146&164&143&178&183&129&165\\
March&100&109&366&228&127&153&246&183\\\hline
Mean&157&229&238&211&156&174&213&183
\end{tabular}
\end{center}

Among the individual observations, the \textit{largest} fluctuations that occurred in the morning and afternoon declinations over the three years may also be given here. The morning declination on 8 October 1835 was $20'1''$ greater than on 7 October, and the afternoon declination on 24 April 1836 was $13'0''$ greater than on the preceding day. On the other hand, the morning or afternoon declinations on two consecutive days have often been exactly equal as well. In the monthly mean fluctuations these extremes naturally come much closer together. Even so, the large differences between the individual months in this respect remain quite remarkable. According to the table above, the mean fluctuation in morning declination in March 1837 was $6'6''$, while in December 1836 it was only $1'11''$.

The results for our two hours do not yet allow us to decide with certainty whether larger fluctuations generally prevail at one time of day than at another. The two means are
\SourceLeaf{573}nearly equal in the first year. In the other two years the morning fluctuations predominate, but the difference between the final results for all three years, $3'31''$ and $3'3''$, is too small to be regarded as established by so few years, although in the means for the individual months in the fourth and eighth columns ten months gave a difference in the same direction.

Combining the morning and afternoon observations gives the following mean fluctuations.
\begin{center}\small
\begin{tabular}{l|r|r|r||r}
&Year I&Year II&Year III&Mean\\\hline
April&108&114&237&164\\
May&176&156&245&196\\
June&139&161&208&172\\
July&173&215&270&223\\
August&224&214&289&244\\
September&167&251&185&204\\
October&152&254&229&216\\
November&133&190&235&191\\
December&160&271&120&195\\
January&160&245&146&189\\
February&150&166&148&155\\
March&114&133&312&206
\end{tabular}
\end{center}
\Needspace{7\baselineskip}
\begin{center}\small
\textit{Mean values.}\par\smallskip
\begin{tabular}{l|r|r|r||r}
&Year I&Year II&Year III&Mean\\\hline
July---December&170&234&228&213\\
Other months&143&167&223&181\\
Whole year&158&204&226&198
\end{tabular}
\end{center}

According to the numbers in the fourth column, the fluctuations are somewhat larger from July to December than in the other six months. The means $3'33''$ and $3'1''$ are probably too close, however, to conclude with certainty that this season is more conducive to larger fluctuations, especially since the difference was pronounced in this way mainly in the one year 1835---1836.

The difference in variability when the three individual years are compared is quite clear, on the other hand. The mean for the third year is nearly half as large again as the mean for the first. The overall
\SourceLeaf{574}mean from all the observations so far, $3'18''$, might therefore still change considerably if the observations were continued for longer.

These are the results that can be drawn from the daily records of magnetic declination so far. It is very desirable that similar work be carried out at several places, and it has already begun recently at some of them. If the observations are made exactly at the same time as ours instead of by local time, as is done in Milan, the comparison of individual days offers opportunities for other combinations, which. once they can cover a fairly considerable length of time, will be of great interest. Observers who follow a similar practice, that is, who make their records at times that agree with ours, are therefore asked to report the results of each day separately. It is enough, however, to give them only in scale divisions, so that conversion into divisions of arc can be dispensed with, provided that the necessary reduction elements are noted as well.

